% Capitulo 9: conclusion.
\section{Conclusión}\label{sec:conclusion}

Este documento empezó como un artículo neutro que enseñaba piezas
sueltas y termina como una pieza editorial que las usa todas a la vez. La
portada con su banda granate, el resumen de un vistazo, el índice, las
nueve secciones, las dos tablas, los cuatro diagramas, las ecuaciones con
sus referencias, los seis listados y las seis citas salieron de una sola
orden: \texttt{texforge build}. Nada se pegó como imagen y nada se
compuso fuera del flujo que el propio documento describe.

texforge demuestra aquí lo que promete en su manual: escribir en texto
plano, revisar sin compilar, formatear sin pensar y componer sin una
distribución externa. Los diagramas viajan con el fuente, el código se
resalta solo, la ortografía se revisa en el idioma del documento y el PDF
resultante se inspecciona desde la terminal. Quien publique con estas
piezas --- una tesis, un informe, un manual --- hereda el mismo andamiaje:
una paleta, dos familias tipográficas y un preámbulo que ya resolvió
las decisiones visuales.

Queda trabajo por delante, como en todo proyecto vivo, pero el punto de
partida es este: un documento real, de presentación cuidada, producido
íntegramente con la herramienta que describe. Quien lo compila obtiene lo
mismo en cualquier máquina, porque el motor viaja con la herramienta y el
fuente no depende de nada instalado a mano: el binario y el texto bastan.
Si algo de lo que ves te resulta útil, su fuente está en
\texttt{examples/texforge-capabilities} esperando tu próxima compilación,
y el manual sigue en \texttt{docs/cli-reference.md}.
